\subsection{二交換子の非共有Segre射影}
\begin{proposition}\label{prop:cubic-CY-two-disjoint-Segre-empty}
二交換子を持つCY包絡のうち、対称射影が
$\langle ac,ad,bc,bd\rangle$ となるものはH3を与えない。
\end{proposition}
\begin{proof}
固定出典\cite{AS4CYpinned}の非共有Segre必要被覆により、
$Q_q,N_1,N_2$ の三表示を調べればよい。
$Q_q$ では $ab=ba,cd=dc$ であり、交差交換係数は
$ac=qca$, $ad=q^{-1}da$, $bc=q^{-1}cb$, $bd=qdb$、$q\ne1$ である。
PBW基底より中心三次は、その中の中心単項式の線型包である。
$a^i b^j c^k d^l$ が中心である条件は $q^{i-j}=q^{k-l}=1$。
総次数3では二つの差の一方は奇数である。$q\ne1$ のため
その差の絶対値は3でなければならず、$q$ は原始三乗根となる。
従って中心三次は $\langle a^3,b^3,c^3,d^3\rangle$ に限られる。

この原始三乗根の場合、順序 $(a,b,c,d)$ で六交換指数は
$(0,1,2,2,1,0)$ となる。
定理\ref{thm:skew-finite-37}の最終行に一致するので、命題\ref{prop:third-six-rank-obstruction}により
全切断が排除される。

$N_1,N_2$ では $a,c$ がそれぞれ正常である。規約を
$nv=\nu_n(v)n$ とすると、$\nu_a$ は $(a,b,c,d)$ 上の
$\operatorname{diag}(1,1,-1,-1)$ である。
$\nu_c$ は $a\mapsto-a,c\mapsto c,d\mapsto d$ とし、
$b\mapsto-b$（$N_1$）又は $b\mapsto-b+8a$（$N_2$）である。
よって非零の $ac$ の正常化作用は半単純部分が $-I$ である。
補題\ref{lem:normalizer-minus-unipotent-odd-centre}から両表示に非零中心三次はない。
\end{proof}

\begin{proposition}\label{prop:cubic-CY-two-veronese-one-empty}
二交換子の混合階数一Veronese型にH3包絡は存在しない。
\end{proposition}
\begin{proof}
固定出典\cite{AS4CYpinned}の必要表示VM1を用いる。
$c,d$ の個数で濾過すると、VM1の六関係の先頭は可換関係である。
AS4のHilbert級数との比較により、随伴次数環は $k[a,b,c,d]$ となる。
$a$ は正常で、正常化作用は $c\mapsto c-2b,d\mapsto d+2a$、
$a,b$ は固定される。従って任意の中心元の最高濾過項は
$(-b\partial_c+a\partial_d)$ の核に入る。
この核は $k[a,b,ac+bd]$ である。実際、$k(a,b)$ 上では線型座標変換で
$k(a,b)[ac+bd]$ を得る。各 $c,d$ 次数 $m$ の係数を比較すると、
$(ac+bd)^m$ の両端係数 $a^m,b^m$ は互いに素なので、係数に分母は残らない。

$z=ac+bd$ は $a,b$ と可換である。
上の核の記述に従って中心元の最高項を $a,b,z$ の多項式で差し引けば、
差は引き続き $a$ と可換であり、濾過次数が下がる。
従って中心元は $k[a,b,z]$ に属する。
VM1から $[c,z]=-bz$、$[d,z]=az$ が従う。
従って $[c,a^ib^jz^m]=(2i-j-m)b a^ib^jz^m$。
随伴次数環での代数的独立性により、この作用の核は
$2i=j+m$ の単項式で張られる。
総次数3の条件 $i+j+2m=3$ と合わせると $i=1,j=2,m=0$ だけとなる。
実際 $ab^2$ は $d$ とも可換である。
中心三次切断は従って $ab^2$ の倍数であり、その商は整域でない。
\end{proof}

\begin{proposition}
混合階数二Veronese型の表示VRで与えられる全 $A_t(r_0,r_1,r_2)$ に対し、
中心三次は零である。
\end{proposition}
\begin{proof}
VM1と同じ濾過を用いる。$a$ の正常化作用は
$c\mapsto c+b/(2t),d\mapsto d-a/(2\beta)$ であり、
最高項の不変環は $k[a,b,z]$、$z=tac+\beta bd$ である。
両端係数の互いに素性と濾過降下の議論はそのまま適用できる。
実際に $z$ は $a,b$ と可換であり、後者には
$t/(t-1)=\beta/(2(2\beta-1))$ を使う。
従って中心三次は $f(a,b)+(Aa+Bb)z$ と書ける。

VRによれば、$[c,z]$ の濾過次数一の部分は
$-\frac{t+1}{2t(t-1)}bz$ である。
$[c,a]=-ab/(2t)$ と $[c,b]=b^2/(t-1)$ を合わせると、
$[c,(Aa+Bb)z]$ の濾過次数一の部分は
$(-Aa/(t-1)+Bb/(2t))bz$ である。
$[c,f(a,b)]$ は濾過次数零なので $A=B=0$ が必要となる。
残る単項式 $a^ib^j$ には $i+j=3$ と $i(t-1)=2tj$ が必要となる。
$i=0,3$ は許容域で不可能、$i=1$ は禁止値 $t=-1/3$、
$i=2$ は矛盾となる。従って中心三次は零である。
\end{proof}

\begin{theorem}\label{thm:cubic-CY-two-commutators-empty}
二交換子のCY包絡を持つH3は存在しない。
\end{theorem}
\begin{proof}
固定出典\cite{AS4CYpinned}の二交換子被覆
\texttt{cor:cy-two-commutators-complete}の必要方向により、
既に扱った表示と、混合階数二の孤立表示 $I$ で尽くされる。
$I$ は同出典VTの $B_I,\gamma=-2/9$ を使う。
元16語のトップ行列から独立な六列を取り、
元関係挿入で得た三次・四次商の次元は20,35である。
中心条件の $140\times20$ 行列の20次小行列式の一つは
$-32768/22876792454961$ であり、中心三次は零となる。
この有理数体上の計算は標数零の拡大でも同じである。
公開検証器\\
\texttt{verify\_cubic\_CY\_two\_commutator\_I.py}
で元行列から再現できる。
\end{proof}
